%% LyX 2.4.4 created this file.  For more info, see https://www.lyx.org/.
%% Do not edit unless you really know what you are doing.
\documentclass[english]{article}
\usepackage[T1]{fontenc}
\usepackage[latin9]{inputenc}
\usepackage{enumitem}
\usepackage{amsmath}
\usepackage{amsthm}
\usepackage{amssymb}
\usepackage{cancel}
\usepackage{geometry}
\geometry{verbose,tmargin=1in,bmargin=1in,lmargin=1in,rmargin=1in}

\makeatletter

%%%%%%%%%%%%%%%%%%%%%%%%%%%%%% LyX specific LaTeX commands.
%% Because html converters don't know tabularnewline
\providecommand{\tabularnewline}{\\}

%%%%%%%%%%%%%%%%%%%%%%%%%%%%%% Textclass specific LaTeX commands.
\numberwithin{equation}{section}
\numberwithin{figure}{section}
\newlength{\lyxlabelwidth}      % auxiliary length 

%%%%%%%%%%%%%%%%%%%%%%%%%%%%%% User specified LaTeX commands.
\usepackage{fancyhdr}
\usepackage{lastpage}
\pagestyle{fancy}
\usepackage{enumitem}
\usepackage{ifthen}
%sets math fonts to be more readable
\everymath=\expandafter{\the\everymath\displaystyle}
%\DeclareMathSizes{10}{10}{10}{10}
\usepackage{multicol}

\usepackage{tikz}
\usetikzlibrary{plotmarks}
\usepackage{pgfplots}
\pgfplotsset{compat=newest}

\usepackage{calc}
\usepackage[nomessages]{fp}

\usepackage{advdate}    % Advancing/saving dates
\usepackage{datetime}   % Dates formatting
\usepackage{datenumber} % Counters for dates
% Added by lyx2lyx
\setlength{\parskip}{\smallskipamount}
\setlength{\parindent}{0pt}

\makeatother

\usepackage{babel}
\begin{document}
\renewcommand{\author}{Dr.\,James Ashe}

\newcommand{\classNum}{132}

\newcommand{\classSec}{C and K}

\newcommand{\nameType}{Name:}

\newcommand{\examType}{Midterm}

\renewcommand{\date}{\formatdate{4}{3}{2014}}

%%First page's header%%%%%%%%%%%%%%%%%%%%%%
\fancypagestyle{firststyle} %%header settings for first pagr
{
	\fancyhead[L]{MTH \classNum{}(\classSec{}) \author{}\\
	\textbf{\examType{}} \date{}}
	\fancyhead[C]{\textbf{\nameType}}
	\setlength{\headsep}{14pt}
}
%%%%%%%%%%%%%%%%%%%%%%%%%%%%%%%%%%%%%%%%%%

%%Next page's header%%%%%%%%%%%%%%%%%%%%%%%
\setlist{nolistsep}
\lhead{MTH \classNum{}(\classSec{})} 
\chead{\textbf{\nameType}} 
\cfoot{\thepage{}~of~\pageref{LastPage}} 
\setlength{\headsep}{5pt} 
%%%%%%%%%%%%%%%%%%%%%%%%%%%%%%%%%%%%%%%%%%%

%%Various settings; see the preamble for other settings%%%%%
%%\setenumerate[0]{leftmargin=12pt,itemindent=0pt,labelwidth=0pt}
\renewcommand{\labelenumi}{\textbf{\arabic{enumi}.}}
\renewcommand{\labelenumii}{\textbf{\alph{enumii}.}}
\thispagestyle{firststyle}
%%%%%%%%%%%%%%%%%%%%%%%%%%%%%%%%%%%%%%%%%%
\newcommand{\xmax}{10}
\newcommand{\xmin}{-10}
\newcommand{\xstep}{0} %must be xmin+n
\newcommand{\ymax}{10}
\newcommand{\ymin}{-10}
\newcommand{\ystep}{0} %must be ymin+n

\newcommand{\xspace}{\vspace{1.5cm}}

\emph{You must show work to receive credit. Do not use a calculator/computer
for problems with a }$\bigstar$\emph{. Each problem is of equal value. }

\textbf{Find the slope of the line passing through the given pair
of points .}
\begin{enumerate}[resume]
\item $(1,2)$ and $(3,-4)$\xspace
\end{enumerate}
\textbf{Use the given conditions to write an equation for the line
in slope-intercept form, if possible.}
\begin{enumerate}[resume]
\item Slope$=-2$, passing through $(13,4)$.\vfill
\item Passing through $(-3,0)$ and $(-6,4)$.\vfill
\end{enumerate}
\textbf{Find the equation of the least squares line. }
\begin{enumerate}[resume]
\item \label{enu:The-total-amount}The paired data below consists of the
number or Florida pleasure boats (in tens of thousands) and the number
of yearly watercraft-related manatee deaths from 1995-2000.

\begin{enumerate}[resume]
\item []%
\begin{tabular}{c|c|c|c|c|c|c}
Hours $(x)$ & 68 & 67 & 70 & 73 & 81 & 84\tabularnewline
\hline 
Score $(y)$ & 53 & 35 & 49 & 60 & 67 & 78\tabularnewline
\end{tabular}\xspace
\end{enumerate}
\end{enumerate}
\textbf{Solve the system of two equations in two unknowns.}
\begin{enumerate}[resume]
\item $\begin{array}{r}
x+y=5\\
x-y=10
\end{array}$\vfill
\item $\begin{array}{r}
x+y=2\\
2x+2y=5
\end{array}$\vfill
\end{enumerate}
\newpage

\textbf{Solve the problem.}
\begin{enumerate}[resume]
\item There are 1000 people at a basketball game. The admission prices was
\$20 for non-students and \$10 for students. The admission receipts
were \$17,500. How many students and non-students attended the game?
\vfill
\item Two pounds of rice and three pounds of potatoes cost a total of \$8,
and three pounds of rice and one pound of potatoes cost a total of
\$5. What is the price per pound for rice and the price per pound
of potatoes?\vfill
\item A car manufacturer has decided to come out with a new type of car.
The fixed cost for the production of this new car is \$250,000 and
the variable costs are \$15,000 per car. The company expects to sell
the cars for \$20,000 each. Formulate a function $C(x)$ for the total
cost of producing $x$ new cars and a function $R(x)$ for the total
revenue generated by the sales of $x$ cars. \vfill
\end{enumerate}
\newpage
\begin{enumerate}[resume]
\item After 5 years as an actuary, Sue's salary was \$100,000. After 10
years her salary had risen to \$200,000. Let $y$ represent her salary
after $x$ years on the job. Assuming that a change in her salary
over time can be approximated by a straight line, give an equation
for this line in the form $y=mx+b$.\vfill
\item A package delivery company charges \$5.00 per package to deliver.
The fixed cost to run the delivery truck is \$450 per day. If the
company charges \$7.00 per package, how many packages must be delivered
to make a profit of \$500?\vfill
\item A package delivery company charges \$5 per package to deliver. The
fixed cost to run the delivery truck is \$450 per day. If the company
charges \$7 per package, how many packages must be delivered to break
even?\vfill
\end{enumerate}
\newpage
\begin{enumerate}[resume]
\item Let the supply and demand functions for gasoline be given by 
\[
p=S(q)=\frac{4}{5}q\mbox{ and }p=D(q)=300-\frac{2}{5}q
\]
where $p$ is the price in dollars and $q$ is the number of 42-gallon
barrels. Find the equilibrium quantity and the equilibrium price.\vfill\xspace
\end{enumerate}
\textbf{The sizes of two matrices $A$ and $B$ are given. Find the
sizes of the product $AB$ and the product $BA$, whenever these products
exist.}
\begin{enumerate}[resume]
\item $A$ is $2\times5$, and $B$ is $5\times3$.

\begin{enumerate}[resume]
\item $AB$:\vspace{.5cm}
\item $BA$:\vspace{.5cm}
\end{enumerate}
\end{enumerate}
\textbf{Find the matrix if possible.}
\begin{enumerate}[resume]
\item $\begin{bmatrix}\begin{array}{rr}
3 & 0\\
-2 & 1\\
0 & 3
\end{array}\end{bmatrix}\begin{bmatrix}\begin{array}{rrr}
1 & 0 & 3\\
0 & 1 & 0
\end{array}\end{bmatrix}$\vfill
\item $\begin{bmatrix}\begin{array}{rr}
3 & 0\\
-2 & 1\\
0 & 3
\end{array}\end{bmatrix}\begin{bmatrix}\begin{array}{rr}
1 & 0\\
0 & 0\\
0 & 0
\end{array}\end{bmatrix}$\vfill
\item $\begin{bmatrix}\begin{array}{rrr}
a & b & 1\\
0 & c & 0\\
0 & 1 & d
\end{array}\end{bmatrix}\begin{bmatrix}\begin{array}{rrr}
0 & 0 & 0\\
0 & 0 & 0\\
0 & 0 & 0
\end{array}\end{bmatrix}$ where $a,b,c,d$ are real numbers.\xspace
\end{enumerate}
\newpage

\textbf{Solve the problem.}
\begin{enumerate}[resume]
\item Let $A=\begin{bmatrix}\begin{array}{rr}
-1 & 3\\
2 & 0
\end{array}\end{bmatrix}$ Find $-2A$. \vfill
\item Let $A=\begin{bmatrix}\begin{array}{rr}
-1 & 3\\
2 & 0
\end{array}\end{bmatrix}$ and $B=\begin{bmatrix}\begin{array}{rr}
0 & 2\\
5 & -1
\end{array}\end{bmatrix}$. Find $A+B$. \vfill
\item Let $A=\begin{bmatrix}\begin{array}{rr}
-1 & 3\\
2 & 0
\end{array}\end{bmatrix}$ and $B=\begin{bmatrix}\begin{array}{rr}
0 & 2\\
5 & -1
\end{array}\end{bmatrix}$. Find $2A+B$. \vfill
\end{enumerate}
\textbf{Bonus}
\begin{enumerate}[resume]
\item [$\bigstar$]Find the inverse ($A^{-1}$) of $A=\begin{bmatrix}\begin{array}{rr}
1 & 0\\
2 & 2
\end{array}\end{bmatrix}$ without using a calculator/computer.\vfill
\end{enumerate}

\end{document}
